\section{Introduction}
\label{sec:intro}

Large language model (LLM) serving has become a storage problem. Every request
first runs a \emph{prefill} pass that turns its prompt into key--value (KV)
tensors, and for long prompts this pass dominates cost and time-to-first-token
(TTFT). Serving systems therefore cache KV tensors and reuse them whenever a new
prompt shares a prefix with an old one~\cite{kwon2023vllm,zheng2024sglang}.
\mc~\cite{qin2025mooncake}, the serving platform of Moonshot AI's Kimi and a
FAST'25 Best Paper, pushed this idea to cluster scale: it pools the otherwise idle
CPU DRAM, SSD and RDMA NICs of a GPU cluster into a petabyte-scale,
deduplicated KVCache, and schedules every request around where its reusable
blocks live. On real traces this raised effective request capacity by
59\%--498\% over vLLM baselines~\cite{qin2025mooncake}. The design is no longer
one company's: \mc's store and transfer engine now back cluster-wide KV pools in
vLLM, SGLang and LMCache~\cite{mooncakerepo,liu2025lmcache}.

A shared cache is shared state, and shared state leaks. A line of recent work
shows that prefix caching in a single serving engine creates a TTFT side channel:
an attacker who guesses part of another user's prompt sees a faster first token
when the guess is cached, and can reconstruct system prompts and user inputs
token by token~\cite{song2024earlybird,wu2025promptpeek,zheng2024inputsnatch}.
Defenses followed. Engines now let operators salt cache keys per
tenant~\cite{vllmcachesalt,dynamonvext}; KVGov salts at the point where prompts
diverge~\cite{addagada2026kvgov}; SafeKV shares only blocks it judges
non-sensitive~\cite{chu2025safekv}; CacheSolidarity detects suspicious reuse and
isolates the affected prefixes~\cite{pennas2026cachesolidarity}.

These defenses share two limits. First, each buys confidentiality by giving up
\emph{physical} sharing, either everywhere or wherever a detector fires; in a
\mc-style cluster, where the global pool exists precisely to raise hit rates
beyond what any single node can hold, that cost lands on the system's main
benefit. Second, each reasons about one engine and one channel, TTFT. A
cluster-scale KVCache exposes more. Its scheduler admits or rejects requests
based on predicted TTFT, which depends on the global prefix match. Commercial
APIs report and discount cached tokens in every
response~\needcite{provider prompt-caching docs, e.g.\ OpenAI cached\_tokens
usage field}, and a 2026 vLLM vulnerability turned that count into a
cross-tenant membership oracle~\cite{cve2026105752}. Its blocks live in
tiers whose latencies form a ladder that reflects how recently and how widely a
prefix was used. And its capacity is shared, so one tenant's traffic evicts
another's blocks even when keys are salted.

\paragraph{This paper.} We make two contributions. The first is an attack study of
cluster-scale, tiered KVCache serving (\cref{sec:attacks}). We build six attacks
against a \mc deployment (\A1--\A6), four of which (\A2--\A4, \A6) we have not
found in prior work \todo{re-run the prior-art search before submission}, and show that \tbd{N of 6} succeed with \tbd{AUC $\geq$ X} under
realistic load; the admission and accounting channels require no timing at all,
and the eviction channel survives per-tenant salting.

The second is \sys, a defense built on one observation: \emph{confidentiality
requires isolating what a principal can observe, not what the system stores.}
\sys leaves \mc's physical layer unchanged---blocks remain globally deduplicated,
replicated and transferred over RDMA---and adds an observable layer that makes
every client-visible quantity a function of the principal's own history and a
coarse public load signal. A per-principal Virtual Cache Directory (\vcd)
answers how long a prefix the principal would have hit in a private cache of
fixed quota. A release controller holds the first token until a deadline
computed from that virtual hit length; admission decisions and usage accounting
use it too. Because the physical cache still hits, the GPU finishes early; a
slack-aware scheduler reinvests that time in other requests. To a client, \sys
is indistinguishable from an isolated system; to the operator, it does the
compute of a shared one.

The idea of releasing outputs on a schedule that depends only on public
information has a long history in timing-channel
mitigation~\cite{askarov2010predictive,zhang2011predictive}, and the leakage of
cross-user deduplication was first studied in cloud
storage~\cite{harnik2010dedup}. \sys's contribution is to apply observational
isolation to a disaggregated KV store, where the right schedule comes from a
\emph{counterfactual} cache rather than from padding, where non-timing
observables must be covered by the same counterfactual, and where the withheld
latency can be recovered as throughput.

We implement \sys on vLLM with the \mc Store connector and a \mc-style global
scheduler (\cref{sec:impl}). On \tbd{testbed} with \mc's released traces,
\sys \tbd{brings all six channels to AUC $\leq$ 0.55}, keeps
\tbd{$\geq$90\%} of the prefill GPU time saved by unrestricted global sharing,
and serves \tbd{X\%} more SLO-compliant requests than per-tenant isolation, at a
cost of \tbd{Y\%} higher median TTFT for requests whose only hits are
cross-tenant (\cref{sec:eval}).

\paragraph{Contributions.}
\begin{itemize}
\item The first systematic study of leakage in cluster-scale, tiered KVCache
  serving, including channels through admission control, usage accounting, tier
  latency, replication and eviction (\cref{sec:attacks}).
\item Observational isolation for shared KV caches: a definition relative to a
  public load signal, and a design that meets it while preserving physical
  sharing (\cref{sec:threat,sec:design,sec:security}).
\item \sys, a prototype on vLLM and \mc Store, with a slack-aware scheduler that
  converts withheld latency into capacity (\cref{sec:design,sec:impl}).
\item An evaluation against six attacks and five defenses on \mc production
  traces (\cref{sec:eval}). We will release code, attack harnesses and
  principal-labelled trace derivatives (\cref{app:openscience}).
\end{itemize}
